% Compile with pdfLaTeX. Place member pictures in the same folder if required.

\documentclass[11pt,a4paper]{article}

\usepackage[a4paper,left=2.0cm,right=2.0cm,top=1.55cm,bottom=1.55cm]{geometry}
\usepackage[T1]{fontenc}
\usepackage{newtxtext}
\usepackage{graphicx}
\usepackage{xcolor}
\usepackage{array}
\usepackage{tabularx}
\usepackage{ragged2e}
\usepackage{enumitem}
\usepackage{fancyhdr}
\usepackage{tikz}
\usepackage{setspace}

% -------------------------------------------------
% COLOURS
% -------------------------------------------------
\definecolor{TitleBlue}{RGB}{61,103,154}

% -------------------------------------------------
% GENERAL SETTINGS
% -------------------------------------------------
\setlength{\parindent}{0pt}
\setlength{\parskip}{0pt}
\renewcommand{\arraystretch}{1.0}

% The supplied document uses a clean sans-serif heading style
% and italic serif-style explanatory text.
\newcommand{\blueheading}[1]{%
    {\sffamily\color{TitleBlue}\fontsize{15}{18}\selectfont #1}%
}

\newcommand{\bluebigheading}[1]{%
    {\sffamily\bfseries\color{TitleBlue}\fontsize{16}{19}\selectfont #1}%
}

\newcommand{\instruction}[1]{%
    {\itshape\fontsize{10.5}{12.5}\selectfont #1}%
}

\newcommand{\answerbox}[2][3.0cm]{%
    \vspace{2pt}
    \noindent
    \fbox{%
        \begin{minipage}[t][#1][t]{\dimexpr\textwidth-2\fboxsep-2\fboxrule\relax}
            \vspace{1mm}
            \itshape #2
        \end{minipage}
    }
}

% -------------------------------------------------
% HEADER / FOOTER
% -------------------------------------------------
\pagestyle{fancy}
\fancyhf{}
\renewcommand{\headrulewidth}{0pt}
\renewcommand{\footrulewidth}{0pt}

\fancyhead[L]{%
    \sffamily\bfseries\itshape\fontsize{10.5}{12}\selectfont
    DS with C (TCS-302) \& OOPS with C++ (TCS-307)\\
    DSCPP-III-2026-T147
}
\fancyhead[R]{%
    \sffamily\bfseries\fontsize{10.5}{12}\selectfont
    PROJECT PROPOSAL
}
\fancyfoot[R]{%
    \itshape\fontsize{10}{12}\selectfont
    Page \thepage
}
\setlength{\headheight}{14pt}
\setlength{\headsep}{23pt}
\setlength{\footskip}{24pt}

% -------------------------------------------------
% PAGE 1
% -------------------------------------------------
\begin{document}

\vspace*{4mm}

\bluebigheading{PROJECT AND TEAM INFORMATION}

\vspace{13mm}

\blueheading{Project Title}

\vspace{1mm}

\answerbox[0.5cm]{StudentSetu: A Campus-Wide Skill Barter, Item Hand-off and Peer Tutoring Platform
}

\vspace{13mm}

\blueheading{Student / Team Information}

\vspace{2mm}


% Team information table

% Compact 3-member layout so all members remain on Page 1.

\noindent

\begin{tabularx}{\textwidth}{|>{\RaggedRight\arraybackslash}p{0.69\textwidth}|>{\RaggedRight\arraybackslash}X|}

\hline

\begin{minipage}[t]{\linewidth}

\vspace{1.5mm}

\textit{Team Name:}\\[-1mm]

\textit{Team \#}

\vspace{1.5mm}

\end{minipage}

&

\begin{minipage}[t]{\linewidth}

\vspace{1.5mm}

\textit{GIG}

\vspace{1.5mm}

\end{minipage}

\\

\hline

\begin{minipage}[t][3.25cm][t]{\linewidth}

\vspace{1.5mm}

\textbf{Team member 1 (Team Lead)}\\[-1mm]

\textit{(Parganiha, Gunmay: 2510370542: 2510370542@geu.ac.in, picture):}

\vspace{2.0cm}

\end{minipage}

&

\begin{minipage}[t][3.25cm][t]{\linewidth}

\vspace{1.5mm}

\textit{Parganiha, Gunmay -- 2510370542}\\

\textit{2510370542@geu.ac.in}

\vspace{2mm}

\includegraphics[width=3.0cm,height=1.5cm,keepaspectratio]{Passport Size Photo (1).png}

\end{minipage}

\\

\hline

\begin{minipage}[t][3.25cm][t]{\linewidth}

\vspace{1.5mm}

\textbf{Team member 2}\\[-1mm]

\textit{(Singh Jethi, Ishwar: 2510370458: 2510370458@geu.ac.in, picture):}

\vspace{2.0cm}

\end{minipage}

&

\begin{minipage}[t][3.25cm][t]{\linewidth}

\vspace{1.5mm}

\textit{Singh Jethi, Ishwar -- 2510370458}\\

\textit{2510370458geu.ac.in}

\vspace{2mm}

\includegraphics[width=3.0cm,height=1.5cm,keepaspectratio]{jethi.png}

\end{minipage}

\\

\hline

\begin{minipage}[t][3.25cm][t]{\linewidth}

\vspace{1.5mm}

\textbf{Team member 3}\\[-1mm]

\textit{(Yadav, Gaurav: 25021088: 25021088@geu.ac.in, picture):}

\vspace{2.0cm}

\end{minipage}

&

\begin{minipage}[t][3.25cm][t]{\linewidth}

\vspace{1.5mm}

\textit{Yadav, Gaurav-- 25021088 }\\

\textit{25021088@geu.ac.in}

\vspace{2mm}

\includegraphics[width=3.0cm,height=1.5cm,keepaspectratio]{gaurav.png}

\end{minipage}

\\

\hline

\end{tabularx}

\newpage

% -------------------------------------------------
% PAGE 2
% -------------------------------------------------

\bluebigheading{PROPOSAL DESCRIPTION (10 pts)}

\vspace{12mm}

\blueheading{Motivation (1 pt)}

\vspace{1mm}

\answerbox[7.0cm]{At college, students already help each other out informally all the time. Someone who's strong in DSA but weak in DBMS ends up swapping help with a friend who's the opposite, and nobody pays anyone anything for it. Seniors also end up with a pile of stuff they don't need anymore after a year or two, like lab coats, drafters, old textbooks, calculators, and this stuff would actually be useful to a junior, but it usually just sits unused or gets thrown away because there's no easy way to pass it on. Right now, both of these things only happen if you already know the right person. If you don't have a friend who's good at what you're struggling with, or you don't happen to know a senior with a spare lab coat, you're just out of luck.

\vspace{4mm}


We wanted to solve this properly instead of leaving it up to luck. The problem isn't that students don't want to help each other, it's that there's no system connecting the right people across the whole college, not just inside existing friend groups. This becomes an even bigger problem for anyone new to campus, like first-years, who haven't built those connections yet. We think this is worth solving because it's something almost every student on campus has personally dealt with, and a working solution could genuinely make life easier for people in our own department, not just be some theoretical project we submit and forget.}

\vspace{12mm}

\blueheading{State of the Art / Current solution (1 pt)}

\vspace{1mm}

\answerbox[6.0cm]{There is no dedicated system for it on our campus. Some marketplace apps exist for buying and selling second hand products like OLX,Facebook marketplace,but they are not built for students and they are not built for swapping purpose , they are built for money changing purpose,so skill for skill does not fit into that situation no matter how we will use it. there is also no campus verification anywhere in these apps which means you are trusting some random person of the internet, not exactly comfortable when you are trading with someone you will actually meet up with.


\vspace{2mm}


Peer tutoring platforms exist but are also mostly paid and scheduled, not students just helping each other for free on their own time. we looked around for a bit and honestly,we were not able to find one thing that puts skill bartering, item handoff and tutoring together for a college campus all in one place specifically, everything solves maybe one part, none of them feels like it was actually made for student problems in life so people just keep doing what already exists which is asking around in the friend circle.}

\vspace{12mm}

\blueheading{Project Goals and Milestones (2 pts)}

\vspace{1mm}

\answerbox[5.6 cm]{Our goal is a platform where students can list what they can offer, what they need, or want to give away, and will get matched automatically instead of totally depending on luck or already knowing the right person around them. but it cannot only do direct matches, real trades do not usually work out that easily so we need the system to catch indirect chains also, that is A helping B, B helping C, looping back to A, otherwise a lot of good matches will just get missed.

\vspace{2mm}
Milestones:

Phase 1 - Finalize class design and core data structures (student records, listings, matching graph)

Phase 2 - Build and unit-test modules separately (skill matching, item hand-off, waitlist system)

Phase 2/3 - Integrate everything into one working program with file storage for persistence

Phase 3 - Add the trust layer (email verification and ratings)

Phase 3 - End-to-end testing and final demo.}

\vspace{25mm}

\blueheading{Project Approach (3 pts)}

\vspace{1mm}

\answerbox[9.0cm]{We're building this as one combined C++ project, using OOP for the overall structure and DSA for the actual matching logic underneath. At the top we have a base Listing class, and this gets extended into different listing types: one for skill/tutoring swaps, one for one-way item or notes listings, and one for paid services like tutoring where the price is negotiable. Each type of listing has its own version of matching function, that is the polymorphism part, the system calls the same function no matter what and depending on the listing type different things happen below it, we don't have to code separate matching logic everywhere which will get messy.

\vspace{2mm}

For the matching, students become nodes in a graph and edges represent who can help whom. direct two way matches get checked first, that's the easy case, if there is no direct match we go looking for longer chains, A to B to C then back to A, using depth first search. chain length gets capped because if it did not the search will run way too long once the platform will have a decent number of students on it. item and notes listing works differently, a hash table will handle category search since that needs to be fast, and there is a queue for waitlist so it will be fair when two or three people want the exact same thing.

\vspace{2mm}

We're also adding a basic trust layer: college email verification when signing up, and both people have to confirm before a trade is marked as done. This way the rating system actually means something and people can't fake it. All the data gets saved to files so listings and match history don't disappear every time the program is closed. So basically the whole project has two sides: an OOP side that handles listings and users, and a DSA side underneath that actually does the matching.}

\newpage

% -------------------------------------------------
% PAGE 3
% -------------------------------------------------

\blueheading{System Architecture (High Level Diagram)(2 pts)}

\vspace{1mm}

\answerbox[15.0cm]{
\begin{center}
\includegraphics[
    width=\textwidth,
    height=15.0cm,
    keepaspectratio
]{diagram.png}
\end{center}
}

\vspace{12mm}

\blueheading{Project Outcome / Deliverables (1 pts)}

\vspace{1mm}

\answerbox[6.4cm]{By the end of this project, we'll have a working C++ application called StudentSetu, where a student can make an account, list a skill they can teach, an item they want to give away, or a tutoring request, and get matched with other students automatically. It can be a direct match, or could end up being a multi person chain if there is no direct one available, either way the system should be able to catch it. along with the actual app we are planning to hand in the source code, documentation on the class design and data structures, and something explaining how the matching actually works under the hood, not just what it does but also why it is built that way.

\vspace{2mm}

We want to demo real matches happening, and specifically we want at least one example where the chain matching actually happens, not just the easy direct case, because otherwise how would we know the dfs part will work. along with that demo we will hand in this proposal and the final report covering all the points, what we have decided, how we tested it, and what we found. the end goal is not just something that looks good on paper, we want it to actually run, actually find matches, and if it were live on our own campus, students could use it in day to day life.}

\vspace{30mm}

\blueheading{Assumptions}
\vspace{2.5mm}

\answerbox[2.5cm]{We're assuming that a decent number of students actually use the platform, since the matching system depends on having enough listings to work with. If very few people sign up, there won't be much to match against, and the system won't really show what it can do. We're also assuming that the actual physical hand-off, like meeting up to exchange an item or starting a tutoring session in person, happens outside the app. The software isn't managing or tracking that part, it's only responsible for finding the match and confirming it.}

\vspace{12mm}

\blueheading{References}

\vspace{1mm}

\answerbox[5.0cm]{• Roth, A. E., Sönmez, T., & Ünver, M. U. (2004). Kidney Exchange. NBER Working Paper No. 10002. https://www.nber.org/system/files/working_papers/w10002/w10002.pdf

\vspace{3mm}
• Abraham, D. J., Blum, A., & Sandholm, T. (2007). Clearing Algorithms for Barter Exchange Markets: Enabling Nationwide Kidney Exchanges. Proceedings of the 8th ACM Conference on Electronic Commerce (EC '07). https://arxiv.org/pdf/1605.08863

\vspace{3mm}
• Cormen, T. H., Leiserson, C. E., Rivest, R. L., & Stein, C. Introduction to Algorithms. MIT Press.

\vspace{3mm}
• Zeng, X., Zhou, Y., & Chen, X. (2021). Study on the System Design of Campus Resource Sharing Platform. E3S Web of Conferences, Vol. 235. https://doi.org/10.1051/e3sconf/202123502038}

\end{document}
``
